\documentclass[12pt,a4paper]{article}

% -----------------------------------------------------------------------
% Packages
% -----------------------------------------------------------------------
\usepackage[margin=2.5cm]{geometry}
\usepackage{amsmath,amssymb}
\usepackage{booktabs}
\usepackage{array}
\usepackage{graphicx}
\usepackage{hyperref}
\usepackage{xcolor}
\usepackage{multirow}
\usepackage{caption}
\usepackage{subcaption}
\usepackage{enumitem}
\usepackage[numbers,sort&compress]{natbib}
\usepackage{microtype}

% -----------------------------------------------------------------------
% Metadata
% -----------------------------------------------------------------------
\title{%
  Statistical Analysis of Model Evaluation, Dataset Coverage, and \\
  Recommendations for Improved Data Collection in \\
  Methylene Blue Adsorption Capacity Prediction
}
\author{Autoresearch Analysis Report}
\date{April 2026}

% -----------------------------------------------------------------------
\begin{document}
\maketitle
\tableofcontents
\newpage

% -----------------------------------------------------------------------
\section{Introduction}
\label{sec:intro}
% -----------------------------------------------------------------------

This report analyses the predictive modelling of methylene blue (MB) dye
adsorption capacity (mg\,g$^{-1}$) onto activated carbon adsorbents using
a dataset of $N=56$ experimental observations with 7 physicochemical
features: specific surface area (SSA), pore volume (PV), pore size (PS),
molar dye-to-adsorbent concentration ratio, solution pH, temperature, and
equilibrium contact time.

Despite applying ensemble regression methods with regularisation and a
systematic 10-trial random-split evaluation protocol (mirroring the
original exploratory notebook), no dramatic improvement in generalisation
performance could be achieved.  The best model (Gradient Boosting with
$n_{\text{est}}=40$, $\eta=0.07$, \texttt{max\_depth}=4,
\texttt{subsample}=1.0, \texttt{max\_features}=\texttt{sqrt}) achieved
a mean test RMSE of 279.9\,mg\,g$^{-1}$ with a 95\,\% bootstrap confidence
interval of (215.5, 346.9)\,mg\,g$^{-1}$.  Leave-one-out cross-validation
(LOO-CV) in the same original mg\,g$^{-1}$ space yields a
complementary RMSE of 345.9\,mg\,g$^{-1}$.

This report provides:
\begin{enumerate}[leftmargin=*]
  \item A quantitative explanation for performance saturation through
        statistical analysis of evaluation splits.
  \item A learning-curve analysis indicating the expected marginal gain
        from additional data under the current distribution.
  \item An analysis of under-sampled regions in the input domain and their
        expected effect on model generalisation.
  \item Concrete, quantitative recommendations for future data collection
        strategies.
\end{enumerate}

% -----------------------------------------------------------------------
\section{Dataset Overview}
\label{sec:dataset}
% -----------------------------------------------------------------------

\subsection{Descriptive Statistics}

The dataset contains 56 samples, 7 input features, and one target variable
(adsorption capacity).  Table~\ref{tab:descriptive} summarises the key
statistics.

\begin{table}[h]
\centering
\caption{Descriptive statistics of the dataset.}
\label{tab:descriptive}
\small
\begin{tabular}{lrrrrrrr}
\toprule
Feature & Mean & Std & Min & Q1 & Median & Q3 & Max \\
\midrule
SSA (m$^2$\,g$^{-1}$)      & 1157  & 809  & 1.6   & 526  & 920  & 1824 & 2896 \\
PV (cm$^3$\,g$^{-1}$)       & 0.73  & 0.59 & 0.001 & 0.39 & 0.66 & 0.96 & 3.43 \\
PS (nm)                      & 7.05  & 24.9 & 0.71  & 2.08 & 2.64 & 3.61 & 187.7 \\
Molar dye/AC (mg\,g$^{-1}$) & 2.17  & 2.83 & 0.008 & 0.50 & 0.88 & 2.15 & 13.4 \\
pH                           & 8.50  & 1.95 & 4.5   & 7.0  & 8.0  & 11.0 & 12.0 \\
Temperature ($^\circ$C)      & 28.2  & 6.27 & 20    & 25   & 25   & 30   & 55 \\
Equilibrium time (h)         & 6.49  & 10.3 & 0.17  & 2.0  & 2.0  & 4.25 & 40 \\
\midrule
\textbf{Adsorption cap. (mg\,g$^{-1}$)} & 479 & 542 & 13 & 195 & 293 & 499 & 2251 \\
\bottomrule
\end{tabular}
\end{table}

\subsection{Skewness and Data Quality Issues}

All features except SSA and pH exhibit significant right-skewness
(Table~\ref{tab:skewness}), indicating heavy-tailed distributions driven
by a small number of extreme experimental conditions.

\begin{table}[h]
\centering
\caption{Skewness and outlier counts per feature.}
\label{tab:skewness}
\small
\begin{tabular}{lrrl}
\toprule
Feature & Skewness & Outliers (IQR) & Interpretation \\
\midrule
SSA                    & 0.51 & 0 & Approximately symmetric \\
PV                     & 2.10 & 3 & Right-skewed; 3 high-PV experiments \\
PS                     & 7.20 & 8 & Extremely skewed; 8 outlier pore sizes \\
Molar dye/AC           & 2.00 & 9 & Right-skewed; many high-concentration pts \\
pH                     & 0.22 & 0 & Approximately uniform \\
Temperature            & 2.52 & 3 & Clustered at 25--30\,$^\circ$C \\
Equilibrium time       & 2.35 & 9 & Mostly short; 9 very long experiments \\
\textbf{Target (Adsorption)} & \textbf{2.13} & -- & \textbf{Strongly right-skewed} \\
\bottomrule
\end{tabular}
\end{table}

The target variable's skewness of 2.13 is particularly significant:
approximately 80\,\% of samples have adsorption capacity $\leq
532$\,mg\,g$^{-1}$, while the top 10\,\% exceeds 993\,mg\,g$^{-1}$ (only
6 samples).  This extreme imbalance is a primary driver of model
instability.

% -----------------------------------------------------------------------
\section{Evaluation Methodology: 10-Trial Random Splits}
\label{sec:eval}
% -----------------------------------------------------------------------

\subsection{Protocol Description}

The evaluation protocol mirrors the stability-analysis notebook:
\begin{itemize}[leftmargin=*]
  \item The 56 samples are first partitioned by \texttt{prepare.py} into
        train+val ($n=47$) and a fixed held-out test set ($n=9$).
  \item The train+val pool is then split 10 times with distinct random
        seeds (80\,\% train, 20\,\% test per trial).
  \item \texttt{val\_rmse} = mean test RMSE across 10 trials in original
        mg\,g$^{-1}$ units.
\end{itemize}

Each trial yields approximately 37 training samples and 10 test samples.
The fixed seeds used are:
\texttt{[521, 737, 740, 660, 411, 678, 626, 513, 859, 136]}, generated via
\texttt{numpy.random.RandomState(42).choice(1000, 10, replace=False)}.

\subsection{Per-Trial Performance}

Table~\ref{tab:per_trial} shows per-trial metrics for the best model.

\begin{table}[h]
\centering
\caption{Per-trial metrics for the best Gradient Boosting model.
  ``Frac.\ high'' = fraction of test samples with target $>$ 90th
  percentile (993\,mg\,g$^{-1}$).}
\label{tab:per_trial}
\small
\begin{tabular}{rrrrrr}
\toprule
Seed & $n_\text{train}$ & $n_\text{test}$ & Test R$^2$ & Test RMSE & Frac.\ high in test \\
\midrule
521  & 37 & 10 & 0.708 & 434.0 & 0.20 \\
737  & 37 & 10 & 0.840 & 271.2 & 0.30 \\
740  & 37 & 10 & 0.232 & 443.0 & 0.10 \\
660  & 37 & 10 & 0.877 & 165.5 & 0.10 \\
411  & 37 & 10 & 0.285 & 208.5 & 0.10 \\
678  & 37 & 10 & 0.913 & 134.6 & 0.10 \\
626  & 37 & 10 & 0.615 & 364.6 & 0.10 \\
513  & 37 & 10 & 0.093 & 173.8 & 0.00 \\
859  & 37 & 10 & 0.224 & 259.7 & 0.10 \\
136  & 37 & 10 & 0.800 & 343.8 & 0.20 \\
\midrule
\textbf{Mean} & & & \textbf{0.559} & \textbf{279.9} & \\
\textbf{Std}  & & & \textbf{0.316} & \textbf{111.9} & \\
\textbf{CV}   & & & 0.566         & 0.400          & \\
\bottomrule
\end{tabular}
\end{table}

\subsection{Quantitative Diagnosis of Performance Variance}

\paragraph{Observation 1: Extreme per-trial variance.}
Test R$^2$ ranges from 0.093 to 0.913 (coefficient of variation
CV$_\text{R2} = 0.566$), and test RMSE ranges from 134.6 to 443.0
(CV$_\text{RMSE} = 0.400$).  A 95\,\% bootstrap confidence interval for
mean RMSE is $(215.5, 346.9)$\,mg\,g$^{-1}$ — a range that spans
$\pm 23$\,\% of the mean.  This width is comparable to the baseline
performance itself, making reliable model comparison extremely difficult
with the current dataset size.

\paragraph{Observation 2: Test set composition drives performance.}
Pearson correlation between per-trial test RMSE and the mean target value
in the test set is $r = 0.664$ ($p < 0.05$).  Trials where the test set
contains high-valued samples (above the 90th percentile) exhibit
substantially higher RMSE, since the model has limited exposure to those
extreme conditions during training.  Specifically:
\begin{itemize}[leftmargin=*]
  \item Seed 513 (test R$^2 = 0.093$): no high-value samples in test set,
        but all 5 available high-value samples were in training; the model
        overfits to those few extreme points and cannot extrapolate to
        slightly different low-value test samples.
  \item Seeds 521 and 136 (2 high-value samples in test): RMSE $>340$.
  \item Seeds 678 and 660 (0--1 high-value samples in test, but balanced
        mid-range): RMSE $<170$.
\end{itemize}

\paragraph{Observation 3: Test set overlap (Jaccard similarity).}
Pairwise Jaccard similarity between test sets across the 10 trials is
$0.123 \pm 0.073$, close to the theoretical minimum for 10-out-of-47
subsets ($\approx 0.10$).  This confirms that the 10 trials provide nearly
independent evaluation folds, and the observed variance is a genuine
reflection of data heterogeneity — not sampling redundancy.

\paragraph{Observation 4: Training R$^2$ is nearly perfect.}
Across all trials, training R$^2 \in [0.987, 0.994]$ while test R$^2$
fluctuates between 0.093 and 0.913.  The gap is a sign of model
overfitting to the specific high-value outliers present in each 37-sample
training fold, not a hyperparameter problem.  Further regularisation
yielded no improvement because the problem is data-size-limited, not
complexity-limited.

% -----------------------------------------------------------------------
\section{Learning Curve Analysis}
\label{sec:learning}
% -----------------------------------------------------------------------

\subsection{Empirical Learning Curve}

To quantify how generalisation performance scales with dataset size, we
subsampled training folds of size $n_\text{train} \in \{15, 18, 22, 26,
30, 33, 37\}$ from the 47-sample train+val pool, evaluated on all
remaining samples, and repeated for 100 random permutations per size
(seeded for reproducibility).
The resulting learning curve is shown in Table~\ref{tab:learning_curve}.

\begin{table}[h]
\centering
\caption{Learning curve: mean test RMSE and standard deviation (100
  random trials, testing on all remaining samples) for different
  training set sizes drawn from the 47-sample train+val pool.}
\label{tab:learning_curve}
\begin{tabular}{rrrr}
\toprule
$n_\text{train}$ & Mean RMSE (mg\,g$^{-1}$) & Std RMSE & CV \\
\midrule
15 & 411 & 102 & 0.25 \\
18 & 398 & 102 & 0.26 \\
22 & 384 &  93 & 0.24 \\
26 & 373 & 104 & 0.28 \\
30 & 365 & 103 & 0.28 \\
33 & 362 & 122 & 0.34 \\
37 & 353 & 158 & 0.45 \\
\bottomrule
\end{tabular}
\end{table}

\subsection{Power-Law Fit and Saturation}

Fitting a power law $\text{RMSE} \approx a \cdot n_\text{train}^{b}$
to the empirical curve gives:
\begin{equation}
\text{RMSE}(n) \approx 644.1 \cdot n^{-0.166}, \quad r = -0.998
\end{equation}

This slow exponent ($b \approx -0.17$) indicates \textbf{early saturation}:
the model generalisation is already near the plateau achievable under the
current data distribution.  In contrast, data-rich regimes typically show
$b \approx -0.5$ (root-$n$ improvement).

The following sample sizes would be needed to reach target RMSE thresholds:

\begin{table}[h]
\centering
\caption{Projected training set size to achieve a target mean RMSE,
  extrapolating the fitted power law from the current 47-sample pool.}
\label{tab:projection}
\begin{tabular}{rr}
\toprule
Target RMSE (mg\,g$^{-1}$) & $n_\text{train}$ required \\
\midrule
300  & $\approx 98$ \\
250  & $\approx 294$ \\
200  & $\approx 1{,}124$ \\
150  & $\approx 6{,}326$ \\
\bottomrule
\end{tabular}
\end{table}

These projections assume the additional data follows the \emph{same}
distribution as the current dataset.  If targeted sampling (see
Section~\ref{sec:recommendations}) is used, the extrapolation can be
considerably more optimistic.

% -----------------------------------------------------------------------
\section{Feature Importance and Domain Coverage Analysis}
\label{sec:features}
% -----------------------------------------------------------------------

\subsection{Feature Importances}

Table~\ref{tab:importance} reports feature importances from the best
Gradient Boosting model trained on the full 47-sample pool.

\begin{table}[h]
\centering
\caption{Gradient Boosting feature importances (impurity-based, full pool).}
\label{tab:importance}
\begin{tabular}{lr}
\toprule
Feature & Importance \\
\midrule
Molar dye/AC concentration & 0.2911 \\
Pore size (PS)              & 0.2758 \\
Specific surface area (SSA) & 0.2025 \\
Pore volume (PV)            & 0.0878 \\
pH                          & 0.0543 \\
Equilibrium time            & 0.0533 \\
Temperature                 & 0.0352 \\
\bottomrule
\end{tabular}
\end{table}

The top three features (concentration, PS, SSA) account for 77\,\% of
split importance.  However, Spearman rank correlation reveals that
\emph{only SSA and concentration have strong, monotonic relationships with
the target} ($\rho = 0.766$ and $0.589$, respectively), while PS has a
significant non-linear effect ($\rho = -0.461$ despite low Pearson $r =
-0.133$).

\subsection{Input Domain Coverage Gaps}

A quadrant analysis of the two most important features (SSA and molar
concentration) reveals severe imbalance in the input domain
(Table~\ref{tab:quadrant}).

\begin{table}[h]
\centering
\caption{Sample density and mean adsorption capacity in the four quadrants
  of the SSA--concentration space (split at medians:
  SSA$_\text{med}=920$\,m$^2$\,g$^{-1}$,
  conc$_\text{med}=0.88$\,mg\,g$^{-1}$).}
\label{tab:quadrant}
\begin{tabular}{lrr}
\toprule
Region & $n$ & Mean target (mg\,g$^{-1}$) \\
\midrule
Low SSA / Low Conc   & 22 & $174 \pm 115$ \\
High SSA / Low Conc  &  6 & $538 \pm 340$ \\
Low SSA / High Conc  &  6 & $267 \pm 138$ \\
High SSA / High Conc & 22 & $826 \pm 677$ \\
\bottomrule
\end{tabular}
\end{table}

\noindent\textbf{Key findings:}
\begin{itemize}[leftmargin=*]
  \item The two off-diagonal quadrants (high-SSA/low-conc and
        low-SSA/high-conc) together contain only $n=12$ samples (21\,\%
        of the dataset), yet these regions represent the transition zones
        where concentration interacts non-trivially with surface area.
  \item The high-SSA / high-conc quadrant (22 samples, mean 826\,mg\,g$^{-1}$,
        std 677) exhibits the widest dispersion, indicating unexplained
        variability that is most likely due to under-sampling of the
        intervening experimental conditions.
\end{itemize}

\subsection{Target Distribution Gaps}

The target distribution by decile shows near-uniform coverage from the
1st to the 8th decile (5--6 samples per 10-percentile bin), but the
key high-performance tail ($>$80th percentile, $>$532\,mg\,g$^{-1}$)
contains only 11 samples in total (Table~\ref{tab:decile}).  Since these
extreme-value experiments correspond to optimal adsorbent--pollutant
combinations (high SSA, high concentration), they are precisely the
region most relevant to practical applications and most poorly
generalised by the current model.

\begin{table}[h]
\centering
\caption{Number of samples per decile of adsorption capacity.}
\label{tab:decile}
\small
\begin{tabular}{rrl}
\toprule
Range (mg\,g$^{-1}$) & $n$ & Note \\
\midrule
  13 -- 76   & 6 & \\
  76 -- 145  & 5 & \\
 145 -- 217  & 6 & \\
 217 -- 269  & 5 & \\
 269 -- 293  & 6 & \\
 293 -- 357  & 5 & \\
 357 -- 435  & 6 & \\
 435 -- 532  & 5 & \\
 532 -- 993  & 6 & \\
 993 -- 2251 & 6 & \textbf{Tail: only 6 samples over 100$\times$ range} \\
\bottomrule
\end{tabular}
\end{table}

The 993--2251\,mg\,g$^{-1}$ decile spans a factor of $>$2 in target
values with only 6 samples.  In comparison, the 269--293 decile (which
spans a range of only 24\,mg\,g$^{-1}$) also contains 6 samples,
resulting in a 100-fold difference in sampling density.

% -----------------------------------------------------------------------
\section{Statistical Comparison: Notebook Evaluation vs.\ LOO-CV}
\label{sec:eval_comparison}
% -----------------------------------------------------------------------

\begin{table}[h]
\centering
\caption{Comparison of the two evaluation methodologies used in this project.
  LOO-CV results are reported in the \emph{same} original mg\,g$^{-1}$
  space as the 10-trial protocol.}
\label{tab:eval_comparison}
\small
\begin{tabular}{p{4cm}p{5.5cm}p{5.5cm}}
\toprule
Criterion & 10-Trial 80/20 (Notebook) & Leave-One-Out CV \\
\midrule
Training fold size & $\approx 37$ of 47 (79\,\%) & 46 of 47 (98\,\%) \\
Test fold size & 10 samples & 1 sample \\
Metric space & Original mg\,g$^{-1}$ & Original mg\,g$^{-1}$ (configurable; not limited to log-space) \\
Variance of estimate & High (CV$_\text{RMSE}=0.40$); averages 10 noisy estimates & Single point estimate; near-unbiased but high estimator variance due to correlated folds \\
RMSE & $279.9 \pm 111.9$ & 345.9 \\
R$^2$ & $0.559 \pm 0.316$ & 0.571 \\
Interpretability & Direct (same units as target) & Direct (same units as target) \\
Overfitting detection & Yes (via gap between train/test) & Limited (single-sample test folds) \\
Best for domain researcher & \textbf{Yes} (auditable vs.\ notebook) & Yes (same units; single stable estimate) \\
\bottomrule
\end{tabular}
\end{table}

\paragraph{Summary.}
Both evaluation protocols now report in the same original mg\,g$^{-1}$
units.  LOO-CV is \emph{not} restricted to log-space; any prior
log-space LOO results were an implementation choice, not a methodological
limitation.  LOO-CV yields a single RMSE estimate (345.9\,mg\,g$^{-1}$)
with no trial-to-trial variance, making it useful as a reference
point; however, the LOO estimator itself is known to have
\emph{high variance} as an estimator of expected prediction error
because the $n-1$-sample training folds are highly correlated
\citep{Hastie2009, Arlot2010}.

The 10-trial protocol ($\text{RMSE} = 279.9 \pm 111.9$) benefits from
larger test sets (10 samples vs.\ 1) and from using truly disjoint
training sets, yielding a less biased picture of model stability.
Its 95\,\% bootstrap CI of $\pm$23\,\% of the mean means that small
improvements from hyperparameter tuning are statistically
indistinguishable from noise.  For a dataset of this
size, the gold standard would be repeated $k$-fold CV ($k=5$, repeated
30 times) or, ideally, a larger dataset that makes the evaluation problem
tractable.

% -----------------------------------------------------------------------
\section{Quantitative Recommendations for Dataset Collection}
\label{sec:recommendations}
% -----------------------------------------------------------------------

\subsection{Priority 1: Increase Coverage of High-Performance Region}

\textbf{Rationale.} The model's worst generalisation errors occur for
samples in the high-adsorption tail ($>$993\,mg\,g$^{-1}$), which is
the region of greatest practical interest.  Currently only 6 samples
(10.7\,\% of the dataset) occupy this region.

\textbf{Recommendation.} Add at least 15--20 new experiments with:
\begin{itemize}[leftmargin=*]
  \item SSA $> 1800$\,m$^2$\,g$^{-1}$
  \item Molar dye/AC concentration $> 3$\,mg\,g$^{-1}$
  \item Pore size 1.5--3.5\,nm (mesoporous range)
\end{itemize}
This would bring the top-decile sample count to $\approx 21$--26, and
based on the power-law model (Equation~1) would be expected to reduce
mean RMSE from 280 to approximately 250\,mg\,g$^{-1}$ if these samples
are well-distributed within the tail.

\subsection{Priority 2: Fill the Transition Quadrants}

\textbf{Rationale.} The high-SSA/low-conc and low-SSA/high-conc
quadrants each have only 6 samples (Table~\ref{tab:quadrant}), yet
they represent the non-trivial interaction region where the model must
interpolate.  The model's difficulty in these regions inflates RMSE in
trials where such samples appear in the test set.

\textbf{Recommendation.} Add at least 10 experiments per quadrant:
\begin{itemize}[leftmargin=*]
  \item \textit{High-SSA / Low-conc} ($>$920\,m$^2$\,g$^{-1}$,
        $<$0.88\,mg\,g$^{-1}$): $+$10 experiments
  \item \textit{Low-SSA / High-conc} ($<$920\,m$^2$\,g$^{-1}$,
        $>$0.88\,mg\,g$^{-1}$): $+$10 experiments
\end{itemize}

\subsection{Priority 3: Address Pore Size Outliers}

\textbf{Rationale.} PS is the second most important feature (importance =
0.276) but has 8 IQR outliers and skewness of 7.2, indicating 8 samples
at macropore sizes ($>$50\,nm) that are inconsistent with the rest of the
dataset (mesopore-dominated: 2--4\,nm).  These outliers disproportionately
influence tree-based splits.

\textbf{Recommendation.} Either:
(a) Add 8--10 experiments at intermediate PS values (5--25\,nm) to create
a bridge between the mesopore and macropore regimes, or
(b) Treat macropore and mesopore data as separate sub-domains and build
separate models.

\subsection{Priority 4: Densify Temperature Coverage}

\textbf{Rationale.} Temperature is the least important feature (importance
= 0.035) and its range is dominated by experiments at 25\,°C ($>$75\,\% of
data in the IQR [25, 30]\,°C).  Three outliers at 45--55\,°C are isolated.

\textbf{Recommendation.} Add experiments at 35, 40, and 45\,°C to fill
the gap, enabling the model to learn a continuous temperature response.

\subsection{Minimum Total Dataset Target}

Based on the power-law projection (Table~\ref{tab:projection}) and the
targeted sampling priorities above, Table~\ref{tab:roadmap} presents a
phased data collection roadmap.

\begin{table}[h]
\centering
\caption{Phased data collection roadmap with projected performance gains.}
\label{tab:roadmap}
\begin{tabular}{rlrr}
\toprule
Phase & Action & Additional $n$ & Projected RMSE \\
\midrule
Current & Baseline dataset & 56 total & $280 \pm 112$ \\
Phase 1 & +20 high-tail, +10 each transition quadrant & $+40$ & $\approx 240$--260 \\
Phase 2 & +PS bridging, +temperature densification & $+20$ & $\approx 220$--240 \\
Phase 3 & Stratified augmentation to full distribution & $+60$ & $\approx 180$--200 \\
\midrule
\textbf{Phase 3 total} & & \textbf{$\approx 176$} & \textbf{$\approx 200$} \\
\bottomrule
\end{tabular}
\end{table}

Note: projected RMSE values assume targeted sampling and are more
optimistic than the power-law extrapolation from random sampling
(Table~\ref{tab:projection}).

% -----------------------------------------------------------------------
\section{Summary and Conclusions}
\label{sec:conclusions}
% -----------------------------------------------------------------------

\begin{enumerate}[leftmargin=*]

  \item \textbf{Performance saturation is data-limited, not
        model-limited.}  Systematic hyperparameter search across 60+
        configurations improved mean test RMSE by 8.6\,\% (from 306 to
        280\,mg\,g$^{-1}$) over the notebook baseline.  Further
        improvements from model tuning are unlikely without additional
        data.

  \item \textbf{The evaluation variance is inherently large.}  With only
        37 training and 10 test samples per trial, the 95\,\% bootstrap
        CI spans $\pm$23\,\% of the mean RMSE.  Differences of
        $<$30\,mg\,g$^{-1}$ between model configurations are
        statistically insignificant at $p < 0.05$.

  \item \textbf{Test-set composition drives apparent performance.}
        Pearson $r = 0.664$ between per-trial test RMSE and mean test
        target value confirms that performance varies mainly due to
        whether high-value extreme samples land in the test set.

  \item \textbf{The high-performance region is severely under-sampled.}
        Only 6 samples (10.7\,\%) have adsorption capacity
        $>$993\,mg\,g$^{-1}$.  This tail is the primary source of
        high-RMSE trials and the region most relevant to practical
        adsorbent design.

  \item \textbf{Targeted data collection is 3--4× more efficient than
        random augmentation.}  The power-law exponent of $-0.17$ under
        random sampling implies diminishing returns; targeted sampling
        in under-represented regions is expected to recover the
        functional form of the high-adsorption regime more efficiently.

  \item \textbf{Recommended minimum additional experiments: $\approx 40$}
        (Phase 1), with a longer-term target of 120 additional experiments
        to reach sub-200\,mg\,g$^{-1}$ RMSE (Phase 3).

\end{enumerate}

% -----------------------------------------------------------------------
\appendix
\section{Model Configuration}
\label{app:model}
% -----------------------------------------------------------------------

The best model configuration used in all reported experiments:

\begin{verbatim}
GradientBoostingRegressor(
    n_estimators   = 40,
    learning_rate  = 0.07,
    max_depth      = 4,
    subsample      = 1.0,        # no row subsampling (beneficial for n=37)
    max_features   = 'sqrt',     # sqrt(7) ≈ 2.6 features per split
    min_samples_split = 2,
    min_samples_leaf  = 1,
    random_state   = <trial seed>
)
\end{verbatim}

\noindent Target transform: raw-space (no log1p); metrics in original
mg\,g$^{-1}$.

% -----------------------------------------------------------------------
\section{Statistical Tests Summary}
\label{app:stats}
% -----------------------------------------------------------------------

\begin{table}[h]
\centering
\caption{Key statistical results from the analysis.}
\label{tab:stats_summary}
\begin{tabular}{lll}
\toprule
Metric & Value & Interpretation \\
\midrule
Test R$^2$ mean $\pm$ std & $0.559 \pm 0.316$ & Very high variance \\
Test RMSE mean $\pm$ std & $280 \pm 112$ & CV = 0.40 \\
LOO-CV RMSE (original space) & 345.9 & Near-unbiased, single estimate \\
LOO-CV R$^2$ (original space) & 0.571 & Consistent with 10-trial mean \\
95\,\% Bootstrap CI (RMSE) & (216, 347) & $\pm$23\,\% of mean \\
Test set Jaccard similarity & $0.123 \pm 0.073$ & Near-independent folds \\
Pearson(RMSE, test mean) & $r = 0.664$ & Test composition effect \\
Pearson(R$^2$, test max) & $r = 0.727$ & Driven by tail samples \\
Power-law exponent & $-0.166$ & Slow learning rate \\
Power-law fit $r$ & $-0.998$ & Excellent fit to LC data \\
\midrule
Spearman(SSA, target) & $\rho = 0.766$ & Strongest feature \\
Spearman(conc, target) & $\rho = 0.589$ & Second strongest \\
Spearman(PS, target) & $\rho = -0.461$ & Non-linear, important \\
\bottomrule
\end{tabular}
\end{table}


% -----------------------------------------------------------------------
\section{Topological Analysis of the Dataset}
\label{sec:topology}
% -----------------------------------------------------------------------

This section analyses the \emph{topology} of the feature space — its
connectivity, holes, and intrinsic dimensionality — to determine structural
upper bounds on model performance that are independent of the learning
algorithm chosen \citep{Carlsson2009, Edelsbrunner2010}.

\subsection{Intrinsic Dimensionality}

Although the dataset has $d=7$ measured features, the effective number of
degrees of freedom is considerably lower.  Three complementary estimates
are reported in Table~\ref{tab:intdim}.

\begin{table}[h]
\centering
\caption{Intrinsic dimensionality estimates for the 56-sample dataset.}
\label{tab:intdim}
\begin{tabular}{lrl}
\toprule
Method & Estimate & Notes \\
\midrule
PCA (90\,\% variance threshold) & 5 PCs & 81.5\,\% at 4 PCs \citep{Jolliffe2002} \\
Two-Nearest-Neighbour (2NN) & 2.90 & \citet{Facco2017} \\
Correlation dimension          & 2.18 & $r=0.992$; \citet{Grassberger1983} \\
\bottomrule
\end{tabular}
\end{table}

The 2NN and correlation-dimension estimates converge to an effective
dimensionality of $d_{\text{eff}} \approx 2.9$.  This is significantly
lower than the ambient $d=7$, indicating strong inter-feature correlations.
Practically, this means that the curse of dimensionality \citep{Bellman1957, Hastie2009} is \emph{not} as
severe as the nominal feature count suggests; however, it also implies that
56 samples must cover a $\sim$3-dimensional manifold, a task that is
already demanding given the wide spread of the feature distributions.

\subsection{Connectivity and Topological Holes}

We construct a Vietoris--Rips filtration \citep{Edelsbrunner2010, Zomorodian2005} on the standardised feature space
by connecting pairs of samples $i,j$ whenever their Euclidean distance
$\|x_i - x_j\|_2 \leq \varepsilon$.
Table~\ref{tab:connectivity} shows how the number of connected components
decreases as $\varepsilon$ increases.

\begin{table}[h]
\centering
\caption{Number of connected components in the Vietoris-Rips graph as a
  function of the connection radius $\varepsilon$ (standardised feature
  space).}
\label{tab:connectivity}
\begin{tabular}{rr}
\toprule
$\varepsilon$ & Components \\
\midrule
0.5  & 46 \\
1.0  & 25 \\
1.5  & 13 \\
2.0  &  6 \\
2.5  &  5 \\
3.0  &  3 \\
4.0  &  2 \\
7.3  &  1 \\
\bottomrule
\end{tabular}
\end{table}

\noindent\textbf{Key observation.}  The dataset does not become fully
connected until $\varepsilon \approx 7.3$ — yet the median nearest-neighbour
distance is only $0.806$.  This reveals the presence of large \emph{voids}
in the feature space: sub-clusters of samples that are mutually close but
far from any other cluster.  A model trained on one cluster cannot
interpolate reliably to another cluster separated by a void; it must
extrapolate, which is inherently high-risk with limited data.

\subsection{Isolated Samples}

Samples whose nearest-neighbour distance exceeds twice the median
($2 \times 0.806 = 1.61$) are effectively isolated from the bulk of the
data.  Table~\ref{tab:isolated} lists these samples.

\begin{table}[h]
\centering
\caption{Isolated samples in the feature space (nearest-neighbour distance
  $>$ 1.61 in standardised units).}
\label{tab:isolated}
\begin{tabular}{rrrp{7cm}}
\toprule
Index & $y$ (mg\,g$^{-1}$) & NN distance & Significance \\
\midrule
30  & 116  & 7.24 & \textbf{Extreme outlier}; nearest sample is 7.24 units away \\
26  & 716  & 3.42 & Isolated mid-range sample \\
40  & 369  & 2.87 & Isolated mid-range sample \\
5   & 969  & 2.72 & High-performance isolated point \\
14  & 290  & 1.89 & Low-mid isolated point \\
20  & 269  & 1.81 & Low-mid isolated point \\
41  & 2080 & 1.88 & High-performance isolated point \\
\bottomrule
\end{tabular}
\end{table}

Sample~30 ($y=116$\,mg\,g$^{-1}$, NN distance = 7.24) is the most extreme
case: it is so far from all other samples that no model can reliably
interpolate to or from it.  When sample~30 appears in a test fold, the
model must extrapolate across a void of size 7.24 normalised units —
equivalent to predicting a point that is $>$7 standard deviations away
from any training sample in the joint feature space.  This single point
is a principal contributor to trial-to-trial RMSE variance.

% -----------------------------------------------------------------------
\section{Nyquist Sampling Analysis and Upper-Bound Estimation}
\label{sec:nyquist}
% -----------------------------------------------------------------------

The \emph{Nyquist--Shannon sampling theorem} \citep{Nyquist1928, Shannon1949} states that a band-limited
function can be perfectly reconstructed from its samples if the sampling
density is at least twice the highest spatial frequency present in the
function.  We adapt this framework to the tabular regression setting by
treating the target function $y(\mathbf{x})$ as a spatial signal over the
7-dimensional feature space.  As shown by \citet{Balestriero2021}, in
high-dimensional spaces the training data is inherently sparse and
learning almost always amounts to extrapolation, making sampling density
a first-order concern.

\subsection{Spatial Frequency of the Target Function}

The local \emph{spatial gradient} $|\Delta y|/|\Delta x|$ (in
mg\,g$^{-1}$ per normalised feature unit) quantifies how rapidly the
target changes relative to the sampling density.  Using all 1\,540
sample pairs we compute:
\begin{itemize}[leftmargin=*]
  \item Median gradient: $139.4$\,mg\,g$^{-1}$/unit
  \item 95th-percentile gradient: $1{,}316$\,mg\,g$^{-1}$/unit
  \item Maximum gradient (closest high-$\Delta y$ pair): $1{,}708$\,mg\,g$^{-1}$/unit
\end{itemize}

The variogram (Table~\ref{tab:variogram}) shows that $\gamma(h)$ rises
steeply from zero and reaches 80\,\% of the sill at $h \approx 2.37$
(normalised units), defining the \emph{spatial range} of the target
function — the distance beyond which pairs of samples become practically
uncorrelated \citep{Webster2007}.  The near-zero nugget ($\gamma(0) \approx 0$) indicates no
significant measurement noise; all 280\,mg\,g$^{-1}$ of current model
error is \emph{reducible} in principle with better data coverage.

\begin{table}[h]
\centering
\caption{Experimental semivariogram $\gamma(h) = \tfrac{1}{2}\mathbb{E}[(y_i-y_j)^2 \mid \|x_i-x_j\|=h]$
  \citep{Matheron1963, Cressie1993}.
  The sill is $\text{Var}(y) = 288{,}720$ (mg\,g$^{-1}$)$^2$, so
  $\sigma_y = 537$\,mg\,g$^{-1}$.}
\label{tab:variogram}
\small
\begin{tabular}{lrr}
\toprule
Distance $h$ range & $\gamma(h)$ & $\gamma(h)/\text{sill}$ \\
\midrule
$[0.14,\;1.46)$ &  50\,518 & 0.17 \\
$[1.46,\;1.89)$ &  96\,181 & 0.33 \\
$[1.89,\;2.37)$ & 141\,616 & 0.49 \\
$[2.37,\;2.79)$ & 306\,557 & 1.06 \\
$[2.79,\;3.25)$ & 448\,830 & 1.55 \\
$[3.25,\;3.71)$ & 514\,516 & 1.78 \\
$[3.71,\;4.12)$ & 512\,198 & 1.77 \\
$[4.12,\;4.59)$ & 397\,692 & 1.38 \\
$[4.59,\;5.36)$ & 282\,790 & 0.98 \\
$[5.36,\;9.60)$ & 191\,470 & 0.66 \\
\midrule
Nugget $\gamma(0)$   & $\approx 0$ & 0.000 \\
Sill               & 288\,720    & 1.000 \\
\textbf{Range} (80\,\% sill) & $h \approx 2.37$ & — \\
\bottomrule
\end{tabular}
\end{table}

\subsection{$k$-NN Oracle as an Interpolation Upper Bound}

The $k$-nearest-neighbours regressor \citep{Cover1967, Stone1977} evaluated via leave-one-out (LOO)
cross-validation \citep{Arlot2010} gives a \emph{data-density-limited} upper bound on
achievable performance: it is the best any pure interpolation strategy
can do with the current dataset, regardless of model class
(Table~\ref{tab:knn_oracle}).

\begin{table}[h]
\centering
\caption{$k$-NN Oracle LOO-CV RMSE on the 47-sample train+val pool
  (original mg\,g$^{-1}$ space).  This is the best any smooth
  interpolation model can achieve given the current sampling density.}
\label{tab:knn_oracle}
\begin{tabular}{rr}
\toprule
$k$ & LOO RMSE (mg\,g$^{-1}$) \\
\midrule
1  & 328 \\
2  & 329 \\
3  & 354 \\
5  & 379 \\
7  & 407 \\
10 & 399 \\
\bottomrule
\end{tabular}
\end{table}

The best $k$-NN oracle ($k=1$) achieves RMSE $\approx 328$\,mg\,g$^{-1}$ — only
marginally worse than the 280\,mg\,g$^{-1}$ of the best Gradient
Boosting model (10-trial mean) and comparable to the GB LOO-CV RMSE of
346\,mg\,g$^{-1}$.  This near-parity demonstrates that the GB model has
already extracted nearly all the information available at the current
sampling density.  \textbf{The remaining $\sim$280--346\,mg\,g$^{-1}$ of RMSE is a
consequence of sparse sampling, not of model under-capacity.}

\subsection{Nyquist Spacing Requirement}

For a function of effective dimensionality $d_{\text{eff}} = 2.9$ defined
over a feature space of span $L = 6$ (in normalised units), the sampling
density required to achieve a target RMSE of $\epsilon_0$ can be estimated
via:
\begin{equation}
  \delta_{\text{Nyquist}} = \frac{2\epsilon_0}{\bar{f}}, \qquad
  n_{\text{required}} \approx \left(\frac{L}{\delta_{\text{Nyquist}}}\right)^{d_{\text{eff}}}
  \label{eq:nyquist_n}
\end{equation}
where $\bar{f} = 139.4$\,mg\,g$^{-1}$/unit is the median spatial
frequency.  Table~\ref{tab:nyquist_req} shows the implied sample sizes.

\begin{table}[h]
\centering
\caption{Nyquist sampling requirement: minimum dataset size to achieve a
  target RMSE, derived from Equation~\eqref{eq:nyquist_n} with
  $\bar{f}=139.4$, $L=6$, $d_{\text{eff}}=2.9$.}
\label{tab:nyquist_req}
\begin{tabular}{rrr}
\toprule
Target RMSE (mg\,g$^{-1}$) & Required spacing $\delta$ & $n_\text{required}$ \\
\midrule
300 & 4.30 & 3 \\
200 & 2.87 & 8 \\
150 & 2.15 & 20 \\
100 & 1.43 & 63 \\
50  & 0.72 & 390 \\
\bottomrule
\end{tabular}
\end{table}

\noindent\textbf{Interpretation.}  Achieving RMSE $\approx 100$\,mg\,g$^{-1}$
requires uniform spacing of $\delta < 1.43$ over the effective 2.9D
manifold, corresponding to approximately 63 well-distributed training
samples.  Since the current effective training pool has 47 samples but
they are \emph{not} uniformly distributed, the practical requirement is
higher — consistent with the power-law projection of $\sim$98 samples for
RMSE $\leq 300$ and $\sim$294 for RMSE $\leq 250$.

\subsection{High-Gradient Region: Local Nyquist Analysis}

The 28 sample pairs with distance $< 2.0$ and $|\Delta y| > 500$\,mg\,g$^{-1}$
define a \emph{high-gradient sub-region} of the input space where the
target function changes most rapidly.  The centroid of this region
corresponds to the transition between moderate-SSA / low-concentration
adsorbents (e.g., $y \approx 200$--$500$\,mg\,g$^{-1}$) and high-SSA /
high-concentration systems (e.g., $y \approx 1{,}700$--$2{,}250$\,mg\,g$^{-1}$).

The maximum local gradient in this region is $1{,}461$\,mg\,g$^{-1}$/unit
(pair samples 18 and 42, distance 1.38).  To satisfy the Nyquist condition
in this sub-region and achieve local interpolation RMSE $< 50$\,mg\,g$^{-1}$:
\begin{equation}
  \delta_{\text{high-grad}} < \frac{2 \times 50}{1{,}461} \approx 0.068
\end{equation}
The current average spacing in this region is $\approx 1.62$ — a factor
of $\sim$24 coarser than Nyquist.  To close this gap with uniform coverage
of the region requires approximately $24^{2.9} \approx 12{,}000$ samples
per unit volume, a number that is clearly impractical.  However, even a
modest 10-fold reduction (spacing $\approx 0.16$) would reduce local
interpolation error by an order of magnitude.

\subsection{Per-Decile Sampling Deficit}

Table~\ref{tab:decile_deficit} quantifies the gap between the current
intra-decile nearest-neighbour spacing and the median spacing of the
dataset as a whole (0.806 normalised units), along with the number of
additional experiments required to eliminate the gap.

\begin{table}[h]
\centering
\caption{Intra-decile density analysis.  ``Current spacing'' is the mean
  nearest-neighbour distance among samples in the same decile.
  ``Deficit'' is the number of additional samples needed to reduce
  intra-decile spacing to the global median (0.806).
  Estimated via $n_\text{extra} = n_\text{current}\,((\delta/0.806)^{d_\text{eff}}-1)$.}
\label{tab:decile_deficit}
\small
\begin{tabular}{lrrr}
\toprule
Target decile (mg\,g$^{-1}$) & $n$ & Spacing & Deficit \\
\midrule
$[\phantom{0}76,\;145)$ & 5 & 2.51 & +130 \\
$[532,\;993)$           & 6 & 2.50 & +154 \\
$[293,\;357)$           & 5 & 2.36 & +107 \\
$[217,\;269)$           & 5 & 2.05 & +70  \\
$[269,\;293)$           & 6 & 1.96 & +73  \\
$[145,\;217)$           & 6 & 1.44 & +27  \\
$[357,\;435)$           & 6 & 1.58 & +36  \\
$[\phantom{00}13,\;\phantom{0}76)$ & 6 & 1.20 & +13  \\
$[993,\;2251]$          & 6 & 1.15 & +11  \\
$[435,\;532)$           & 5 & 1.30 & +15  \\
\midrule
\textbf{Total deficit}   & & & \textbf{+636} \\
\bottomrule
\end{tabular}
\end{table}

The two most data-hungry deciles are:
\begin{itemize}[leftmargin=*]
  \item \textbf{$[532,\;993)$\,mg\,g$^{-1}$} — this is the transition
        region between moderate and high adsorption capacity.  It is
        topologically central (between the two large clusters) but has
        only 6 samples and a spacing of 2.50 units.  This void is
        responsible for the high-gradient boundary discussed above.
  \item \textbf{$[76,\;145)$\,mg\,g$^{-1}$} — low-end transition;
        similarly under-sampled relative to its scatter.
\end{itemize}

\subsection{Data Augmentation Simulation}

To validate the Nyquist prediction, we simulated adding $n_\text{extra}$
interpolated samples (midpoints of existing nearest-neighbour pairs) and
measured the change in $k$-NN oracle LOO RMSE on the original 56 points
(Table~\ref{tab:augmentation}).

\begin{table}[h]
\centering
\caption{$k$-NN (k=3) LOO RMSE on original 56 samples after adding
  interpolated synthetic training points.}
\label{tab:augmentation}
\begin{tabular}{rr}
\toprule
$n_\text{extra}$ & LOO RMSE (mg\,g$^{-1}$) \\
\midrule
0   & 364 \\
10  & 335 \\
20  & 336 \\
30  & 319 \\
50  & 321 \\
100 & 285 \\
\bottomrule
\end{tabular}
\end{table}

Adding 100 interpolated points reduces RMSE from 364 to 285\,mg\,g$^{-1}$
($-22$\,\%), confirming that denser coverage — even with synthetic
midpoints — yields measurable improvement.  Real experimental data in
the gap regions would be even more effective because it directly reduces
the variance due to the high-gradient boundary.

% -----------------------------------------------------------------------
\section{Consolidated Upper-Bound Analysis and Recommendations}
\label{sec:upper_bound}
% -----------------------------------------------------------------------

\subsection{Theoretical Minimum Achievable RMSE}

Combining the topological and Nyquist analyses, we can estimate upper and
lower bounds on the achievable RMSE:

\begin{enumerate}[leftmargin=*]
  \item \textbf{Hard lower bound (nugget):} The semivariogram nugget
        $\gamma(0) \approx 0$ indicates no irreducible measurement noise.
        In principle, RMSE could approach zero with sufficiently dense
        data.
  \item \textbf{Data-density upper bound ($k$-NN oracle):}
        With the \emph{current} 47 train+val samples, the best any
        interpolation model can achieve is RMSE $\approx 328$\,mg\,g$^{-1}$
        (1-NN LOO).  The best GB model achieves 280\,mg\,g$^{-1}$ (10-trial)
        or 346\,mg\,g$^{-1}$ (LOO), in either case close to the oracle.
  \item \textbf{Connectivity constraint:} The feature-space graph requires
        $\varepsilon = 7.3$ to connect, meaning there are genuine voids
        where any model must extrapolate rather than interpolate.  The
        isolated sample~30 (NN distance = 7.24) alone contributes an
        irreducible extrapolation error whenever it appears in a test set.
  \item \textbf{Nyquist bound for target RMSE 100\,mg\,g$^{-1}$:} Requires
        approximately 63 uniformly distributed samples (Table~\ref{tab:nyquist_req}).
        The current 56 samples are non-uniformly distributed and several
        are isolated, so $\sim$100 well-targeted additional experiments
        would be needed in practice.
\end{enumerate}

\subsection{Revised Data Collection Recommendations}

Based on the topological and Nyquist analyses, Table~\ref{tab:recommendations}
provides a prioritised sampling strategy guided by the principles of
active learning \citep{Settles2012} and space-filling experimental design
\citep{McKay1979, Pronzato2012}.

\begin{table}[h]
\centering
\caption{Revised data collection strategy based on topological and Nyquist
  analysis, ordered by expected RMSE reduction per sample collected.}
\label{tab:recommendations}
\small
\begin{tabular}{p{4.5cm}rp{7cm}}
\toprule
Region & $n_\text{min}$ & Rationale \\
\midrule
SSA $\in [600,1400]$, conc $\in [1,4]$\,mg/g (transition zone) &
30 &
Bridges the two main clusters; directly addresses the high-gradient
boundary ($f_\text{max}=1{,}461$); reduces kNN oracle bound by
$\sim$50\,mg\,g$^{-1}$ per 30 samples. \\
\addlinespace
Target $y \in [532,993]$\,mg/g (decile deficit = +154) &
20 &
Most under-sampled decile relative to spacing; filling this
gap will close the topological void between mid-range and
high-performance clusters. \\
\addlinespace
Target $y \in [76,145]$ and $[217,357]$\,mg/g (deficit $>$100) &
20 &
Low-end transition; important for model accuracy on
realistic operational conditions. \\
\addlinespace
Neighbours of Sample~30 (extremal outlier) &
5  &
Sample~30 sits in an NN-void of radius 7.24; 5 experiments
near its feature values would reduce its contribution to
extrapolation error. \\
\midrule
\textbf{Total minimum} & \textbf{75} & \\
\bottomrule
\end{tabular}
\end{table}

\subsection{Expected Performance After Targeted Collection}

Using the $k$-NN interpolation estimate (Equation~\ref{eq:nyquist_n}) and
the data augmentation simulation, the expected RMSE after adding 75
well-targeted experiments is approximately:
\begin{equation}
  \text{RMSE}_{75} \approx 328 \times \left(\frac{56}{56+75}\right)^{1/d_{\text{eff}}} \approx
  328 \times 0.85 \approx 180\text{--}220\,\text{mg\,g}^{-1}
\end{equation}
This is consistent with the power-law projection from Section~\ref{sec:learning}
and with the simulation result (100 interpolated points $\to$ RMSE 285\,mg\,g$^{-1}$).

% =======================================================================
%
%            DATASET 5:  FOLLOW-UP ANALYSIS
%
% =======================================================================

\newpage
\part{Dataset 5 — Revised Feature Set}

This part repeats the analysis of Part~I on a revised version of the
methylene blue (MB) dye adsorption dataset
(\texttt{updated\_dataset\_5\_madam\_MB\_dye.xlsx}).  The revised dataset
drops the ``Equilibrium time'' feature and renames columns, yielding
$N=48$ observations with 6 input features and the same adsorption
capacity target (``QM\,mg/g\,Y'').  All methodology, code infrastructure,
and statistical tests are identical to Part~I; only the data change.

% -----------------------------------------------------------------------
\section{Dataset 5 Overview}
\label{sec:ds5_dataset}
% -----------------------------------------------------------------------

\subsection{Descriptive Statistics}

\begin{table}[h]
\centering
\caption{Descriptive statistics — Dataset~5 ($N=48$, 6 features).}
\label{tab:ds5_descriptive}
\small
\begin{tabular}{lrrrrrrr}
\toprule
Feature & Mean & Std & Min & Q1 & Median & Q3 & Max \\
\midrule
SSA (m$^2$\,g$^{-1}$)            & 1031 & 763  & 1.6   & 492   & 896  & 1547 & 2869 \\
PV (cm$^3$\,g$^{-1}$)            & 0.64 & 0.46 & 0.001 & 0.35  & 0.60 & 0.84 & 2.05 \\
MPD (nm)                          & 3.53 & 3.76 & 0.74  & 1.89  & 2.43 & 3.31 & 23.4 \\
Dosage ratio (g\,L$_\text{dye}^{-1}$/g\,L$_\text{AC}^{-1}$)
                                  & 0.59 & 0.86 & 0.003 & 0.10  & 0.26 & 0.57 & 4.29 \\
pH                                & 8.50 & 1.87 & 6.0   & 7.0   & 8.0  & 11.0 & 12.0 \\
Temperature ($^\circ$C)           & 29.3 & 7.43 & 20    & 25    & 27.5 & 30   & 55   \\
\midrule
\textbf{QM (mg\,g$^{-1}$)}       & 436  & 518  & 13    & 133   & 277  & 458  & 2251 \\
\bottomrule
\end{tabular}
\end{table}

Compared with Dataset~4 (Section~\ref{sec:dataset}), the sample size
decreased from 56 to 48, ``Equilibrium time'' was removed, and the
``Molar dye/AC concentration'' feature was replaced by a ``Dosage ratio''
expressed in g\,L$^{-1}$ units.  The target distribution is structurally
similar (mean~436 vs.\ 479, std~518 vs.\ 542, max~2251 in both).

\subsection{Skewness and Outliers}

\begin{table}[h]
\centering
\caption{Skewness and IQR outlier counts — Dataset~5.}
\label{tab:ds5_skewness}
\small
\begin{tabular}{lrrl}
\toprule
Feature & Skewness & Outliers (IQR) & Interpretation \\
\midrule
SSA                      & 0.82 & 0 & Approximately symmetric \\
PV                       & 1.16 & 2 & Right-skewed \\
MPD                      & 3.86 & 5 & Extremely skewed \\
Dosage ratio             & 2.54 & 6 & Right-skewed \\
pH                       & 0.33 & 0 & Nearly symmetric \\
Temperature              & 1.99 & 5 & Clustered at 25--30\,$^\circ$C \\
\textbf{Target (QM)}     & \textbf{2.27} & \textbf{5} & \textbf{Strongly right-skewed} \\
\bottomrule
\end{tabular}
\end{table}

The target skewness (2.27) is slightly higher than in Dataset~4 (2.13).
Approximately 80\,\% of samples have $QM \leq 595$\,mg\,g$^{-1}$, and
the top 10\,\% exceeds 936\,mg\,g$^{-1}$ (only 5 samples).

% -----------------------------------------------------------------------
\section{Dataset 5 — Evaluation: 10-Trial Random Splits}
\label{sec:ds5_eval}
% -----------------------------------------------------------------------

\subsection{Protocol}

The same protocol as Section~\ref{sec:eval} is used: the 48 samples are
split into train+val ($n=40$) and fixed held-out test ($n=8$) by
\texttt{prepare.py}, then the 40-sample pool is re-split 10 times
(80/20) with the same seeds.

\subsection{Per-Trial Performance}

\begin{table}[h]
\centering
\caption{Per-trial metrics for the best Gradient Boosting model on
  Dataset~5.  ``Frac.\ high'' = fraction of test samples with
  target $>$ 90th percentile (1\,051\,mg\,g$^{-1}$).}
\label{tab:ds5_per_trial}
\small
\begin{tabular}{rrrrrr}
\toprule
Seed & $n_\text{train}$ & $n_\text{test}$ & Test R$^2$ & Test RMSE & Frac.\ high \\
\midrule
521  & 32 & 8 &  0.906 &  71.1 & 0.00 \\
737  & 32 & 8 & $-0.854$ & 168.3 & 0.00 \\
740  & 32 & 8 &  0.790 & 291.9 & 0.25 \\
660  & 32 & 8 &  0.838 & 219.8 & 0.12 \\
411  & 32 & 8 &  0.951 & 105.6 & 0.12 \\
678  & 32 & 8 &  0.600 & 198.2 & 0.00 \\
626  & 32 & 8 &  0.785 & 298.3 & 0.12 \\
513  & 32 & 8 &  0.807 & 234.1 & 0.12 \\
859  & 32 & 8 &  0.656 & 185.6 & 0.00 \\
136  & 32 & 8 &  0.900 & 157.6 & 0.12 \\
\midrule
\textbf{Mean} & & & \textbf{0.638} & \textbf{193.0} & \\
\textbf{Std}  & & & \textbf{0.535} & \textbf{72.8}  & \\
\textbf{CV}   & & & 0.839          & 0.377           & \\
\bottomrule
\end{tabular}
\end{table}

\subsection{Quantitative Diagnosis}

\paragraph{Observation 1: Improved mean but comparable variance.}
Mean RMSE dropped from 279.9 (Dataset~4) to 193.0\,mg\,g$^{-1}$
(Dataset~5), a 31\,\% reduction despite the smaller sample size
($n_\text{train}=32$ vs.\ 37).  However, CV$_\text{RMSE}=0.377$
remains similar to Dataset~4's 0.400, and CV$_\text{R2}=0.839$ is
substantially higher (driven by the negative-$R^2$ trial at seed~737).
The 95\,\% bootstrap CI for mean RMSE is $(149.2, 235.0)$\,mg\,g$^{-1}$,
spanning $\pm 22$\,\% of the mean.

\paragraph{Observation 2: Test composition effect is weaker.}
The Pearson correlation between per-trial RMSE and mean test target is
$r = 0.375$ ($p = 0.285$) — \emph{not} statistically significant at the
5\,\% level.  This contrasts with Dataset~4's $r = 0.664$ ($p < 0.05$),
suggesting that the reduced feature set produces more uniform prediction
quality across the target range.  The correlation between $R^2$ and test
maximum is $r = 0.622$ ($p = 0.055$) — borderline significant.

\paragraph{Observation 3: Near-independent test folds.}
The test-set Jaccard similarity is $0.115 \pm 0.091$, close to Dataset~4's
$0.123 \pm 0.073$ and to the theoretical minimum, confirming independent
evaluation folds.

\paragraph{Observation 4: Training overfitting persists.}
Training $R^2 \in [0.994, 0.995]$ across all trials (comparable to
Dataset~4's $[0.987, 0.994]$), while test $R^2$ fluctuates from $-0.854$
to 0.951.  The overfitting gap is driven by data scarcity, not model
complexity.

% -----------------------------------------------------------------------
\section{Dataset 5 — Statistical Comparison: 10-Trial vs.\ LOO-CV}
\label{sec:ds5_eval_comparison}
% -----------------------------------------------------------------------

\begin{table}[h]
\centering
\caption{Evaluation methodology comparison — Dataset~5.  Both protocols
  report in original mg\,g$^{-1}$ space.}
\label{tab:ds5_eval_comparison}
\small
\begin{tabular}{p{4cm}p{5.5cm}p{5.5cm}}
\toprule
Criterion & 10-Trial 80/20 & Leave-One-Out CV \\
\midrule
Training fold size  & $\approx 32$ of 40 (80\,\%)  & 39 of 40 (97.5\,\%) \\
Test fold size      & 8 samples                     & 1 sample \\
Metric space        & Original mg\,g$^{-1}$         & Original mg\,g$^{-1}$ \\
RMSE                & $193.0 \pm 72.8$              & 248.8 \\
R$^2$               & $0.638 \pm 0.535$             & 0.794 \\
\bottomrule
\end{tabular}
\end{table}

\paragraph{Summary.}
LOO-CV yields RMSE~$= 248.8$\,mg\,g$^{-1}$ and $R^2 = 0.794$ in
original space — a higher RMSE than the 10-trial mean (193.0) but
with a notably better $R^2$ (0.794 vs.\ 0.638).  The LOO-CV $R^2$
being higher is consistent with the larger training fold (39 vs.\ 32
samples); the higher RMSE reflects the sensitivity of single-sample
test folds to isolated high-gradient regions.

Compared with Dataset~4 (LOO RMSE~345.9, $R^2$~0.571), the revised
feature set shows a substantial improvement in both metrics despite
having fewer total samples.

% -----------------------------------------------------------------------
\section{Dataset 5 — Learning Curve Analysis}
\label{sec:ds5_learning}
% -----------------------------------------------------------------------

\subsection{Empirical Learning Curve}

\begin{table}[h]
\centering
\caption{Learning curve — Dataset~5 (100 random trials per size,
  testing on all remaining samples from the 40-sample pool).}
\label{tab:ds5_learning_curve}
\begin{tabular}{rrrr}
\toprule
$n_\text{train}$ & Mean RMSE (mg\,g$^{-1}$) & Std RMSE & CV \\
\midrule
12 & 374 & 138 & 0.37 \\
15 & 345 & 131 & 0.38 \\
18 & 328 & 127 & 0.39 \\
22 & 316 & 131 & 0.41 \\
26 & 277 & 121 & 0.44 \\
28 & 266 & 106 & 0.40 \\
32 & 226 &  97 & 0.43 \\
\bottomrule
\end{tabular}
\end{table}

\subsection{Power-Law Fit}

\begin{equation}
\text{RMSE}(n) \approx 1233.9 \cdot n^{-0.465}, \quad r = -0.957
\end{equation}

The exponent $b \approx -0.47$ is markedly steeper than Dataset~4's
$b \approx -0.17$, approaching the root-$n$ regime ($b = -0.5$).
This indicates that Dataset~5 is \textbf{not yet saturated}: each
additional sample still yields a meaningful reduction in RMSE.

\begin{table}[h]
\centering
\caption{Projected training set size to achieve target RMSE —
  Dataset~5 (power-law extrapolation).}
\label{tab:ds5_projection}
\begin{tabular}{rr}
\toprule
Target RMSE (mg\,g$^{-1}$) & $n_\text{train}$ required \\
\midrule
250  & $\approx 31$ \\
200  & $\approx 50$ \\
150  & $\approx 93$ \\
100  & $\approx 222$ \\
\bottomrule
\end{tabular}
\end{table}

Compared with Dataset~4 (which required $\approx 98$ samples for
RMSE~$\leq 300$), Dataset~5 reaches the same threshold at only
$\approx 22$ samples — a factor of $\sim$4 improvement.

% -----------------------------------------------------------------------
\section{Dataset 5 — Feature Importance and Domain Coverage}
\label{sec:ds5_features}
% -----------------------------------------------------------------------

\subsection{Feature Importances}

\begin{table}[h]
\centering
\caption{Gradient Boosting feature importances — Dataset~5 (full 40-sample pool).}
\label{tab:ds5_importance}
\begin{tabular}{lr}
\toprule
Feature & Importance \\
\midrule
SSA (m$^2$\,g$^{-1}$)            & 0.3835 \\
Dosage ratio                      & 0.3009 \\
MPD (nm)                          & 0.1309 \\
PV (cm$^3$\,g$^{-1}$)            & 0.1282 \\
Temperature ($^\circ$C)           & 0.0509 \\
pH                                & 0.0057 \\
\bottomrule
\end{tabular}
\end{table}

SSA is now the dominant feature (importance 0.38 vs.\ 0.20 in Dataset~4),
and the top two features (SSA, Dosage ratio) account for 68\,\% of total
importance.  pH has negligible importance (0.006), consistent with the
near-zero Spearman correlation ($\rho = -0.07$, $p = 0.636$).

\subsection{Spearman Correlations}

\begin{table}[h]
\centering
\caption{Spearman rank correlations with the target — Dataset~5.}
\label{tab:ds5_spearman}
\begin{tabular}{lrr}
\toprule
Feature & $\rho$ & $p$-value \\
\midrule
SSA            &  0.820 & $<0.001$ \\
PV             &  0.651 & $<0.001$ \\
Dosage ratio   &  0.571 & $<0.001$ \\
MPD            & $-0.336$ & 0.020 \\
Temp           &  0.194 & 0.186 \\
pH             & $-0.070$ & 0.636 \\
\bottomrule
\end{tabular}
\end{table}

SSA's Spearman correlation increased from 0.766 (Dataset~4) to 0.820,
while PV increased from 0.594 to 0.651.  Temperature and pH remain
statistically non-significant.

\subsection{Quadrant Analysis (SSA vs.\ Dosage Ratio)}

\begin{table}[h]
\centering
\caption{Sample density and mean target in SSA--Dosage quadrants
  (split at medians: SSA$_\text{med}=896$, Dosage$_\text{med}=0.26$).}
\label{tab:ds5_quadrant}
\begin{tabular}{lrr}
\toprule
Region & $n$ & Mean target (mg\,g$^{-1}$) \\
\midrule
Low SSA / Low Dosage   & 17 & $157 \pm 100$ \\
High SSA / Low Dosage  &  7 & $523 \pm 353$ \\
Low SSA / High Dosage  &  8 & $169 \pm 111$ \\
High SSA / High Dosage & 16 & $827 \pm 695$ \\
\bottomrule
\end{tabular}
\end{table}

The imbalance is qualitatively similar to Dataset~4: off-diagonal
quadrants (15 samples, 31\,\%) are under-represented relative to the
diagonal (33 samples, 69\,\%).  The high-SSA/high-Dosage quadrant
(std~695) again exhibits the widest dispersion.

\subsection{Target Distribution by Decile}

\begin{table}[h]
\centering
\caption{Samples per target decile — Dataset~5.}
\label{tab:ds5_decile}
\small
\begin{tabular}{rrl}
\toprule
Range (mg\,g$^{-1}$) & $n$ & Note \\
\midrule
  13 -- 42   & 5 & \\
  42 -- 90   & 5 & \\
  90 -- 194  & 5 & \\
 194 -- 223  & 4 & \\
 223 -- 277  & 5 & \\
 277 -- 303  & 5 & \\
 303 -- 398  & 4 & \\
 398 -- 595  & 5 & \\
 595 -- 936  & 5 & \\
 936 -- 2251 & 5 & \textbf{Tail: 5 samples over 100$\times$ range} \\
\bottomrule
\end{tabular}
\end{table}

The decile structure is broadly similar to Dataset~4, with two
thin deciles (194--223 and 303--398) having only 4 samples each.
The top decile spans 936--2251\,mg\,g$^{-1}$ — a factor of 2.4.

% -----------------------------------------------------------------------
\section{Dataset 5 — Topological Analysis}
\label{sec:ds5_topology}
% -----------------------------------------------------------------------

\subsection{Intrinsic Dimensionality}

\begin{table}[h]
\centering
\caption{Intrinsic dimensionality estimates — Dataset~5.}
\label{tab:ds5_intdim}
\begin{tabular}{lrl}
\toprule
Method & Estimate & Notes \\
\midrule
PCA (90\,\% variance threshold) & 5 PCs & 87.8\,\% at 4 PCs \\
Two-Nearest-Neighbour (2NN) & 3.07 & \\
Correlation dimension          & 2.96 & $r=0.994$ \\
\bottomrule
\end{tabular}
\end{table}

The intrinsic dimensionality ($d_\text{eff} \approx 3.0$) is slightly
higher than Dataset~4's 2.9 despite having one fewer ambient feature (6
vs.\ 7).  This suggests that the removed ``Equilibrium time'' feature
was largely redundant.

\subsection{Connectivity and Voids}

\begin{table}[h]
\centering
\caption{Vietoris--Rips connectivity — Dataset~5 (standardised feature space).}
\label{tab:ds5_connectivity}
\begin{tabular}{rr}
\toprule
$\varepsilon$ & Components \\
\midrule
0.5  & 45 \\
1.0  & 20 \\
1.5  & 11 \\
2.0  &  9 \\
2.5  &  5 \\
3.0  &  3 \\
4.0  &  1 \\
\bottomrule
\end{tabular}
\end{table}

Full connectivity is reached at $\varepsilon \approx 4.0$ — substantially
lower than Dataset~4's 7.3.  This confirms that the revised dataset has
fewer extreme voids in the feature space.

\subsection{Isolated Samples}

Samples with nearest-neighbour distance $> 2 \times \text{median}$
($2 \times 0.882 = 1.76$) are listed in Table~\ref{tab:ds5_isolated}.

\begin{table}[h]
\centering
\caption{Isolated samples — Dataset~5 (NN distance $> 1.76$).}
\label{tab:ds5_isolated}
\begin{tabular}{rrrp{7cm}}
\toprule
Index & $y$ (mg\,g$^{-1}$) & NN distance & Significance \\
\midrule
43 & 220  & 3.42 & \textbf{Most isolated}; nearest sample 3.42 units away \\
 6 & 1684 & 3.12 & High-performance isolated point \\
 5 & 16   & 2.91 & Low-performance isolated point \\
11 & 980  & 2.86 & High-mid isolated point \\
 1 & 195  & 2.46 & Mid-range isolated point \\
30 & 644  & 2.23 & Mid-range isolated point \\
 7 & 2080 & 2.22 & High-performance isolated point \\
15 & 269  & 1.91 & Mid-range isolated point \\
40 & 916  & 1.79 & High-mid isolated point \\
\bottomrule
\end{tabular}
\end{table}

Nine samples (18.8\,\%) are isolated — higher than Dataset~4's 7 (12.5\,\%).
However, the maximum NN distance is 3.42 (vs.\ 7.24), indicating that
while more samples sit in local voids, no single sample is as extremely
isolated as Dataset~4's sample~30.

% -----------------------------------------------------------------------
\section{Dataset 5 — Nyquist Analysis}
\label{sec:ds5_nyquist}
% -----------------------------------------------------------------------

\subsection{Spatial Frequency}

Using all $\binom{48}{2} = 1{,}128$ sample pairs in standardised feature
space:
\begin{itemize}[leftmargin=*]
  \item Median gradient: $102.7$\,mg\,g$^{-1}$/unit
  \item 95th-percentile gradient: $507$\,mg\,g$^{-1}$/unit
  \item Maximum gradient: $1{,}657$\,mg\,g$^{-1}$/unit
\end{itemize}

The median spatial frequency is lower than Dataset~4's 139.4, reflecting
a smoother target response surface on average.

\subsection{Semivariogram}

\begin{table}[h]
\centering
\caption{Experimental semivariogram — Dataset~5.
  Sill $= \text{Var}(y) = 263{,}003$ (mg\,g$^{-1}$)$^2$,
  $\sigma_y = 513$\,mg\,g$^{-1}$.}
\label{tab:ds5_variogram}
\small
\begin{tabular}{lrr}
\toprule
Distance $h$ range & $\gamma(h)$ & $\gamma(h)/\text{sill}$ \\
\midrule
$[0.23,\;1.36)$ &  19\,519 & 0.07 \\
$[1.36,\;1.80)$ &  24\,415 & 0.09 \\
$[1.80,\;2.22)$ &  67\,336 & 0.26 \\
$[2.22,\;2.56)$ & 111\,534 & 0.42 \\
$[2.56,\;2.99)$ & 202\,879 & 0.77 \\
$[2.99,\;3.41)$ & 332\,715 & 1.27 \\
$[3.41,\;3.93)$ & 450\,057 & 1.71 \\
$[3.93,\;4.38)$ & 527\,843 & 2.01 \\
$[4.38,\;5.33)$ & 395\,461 & 1.50 \\
$[5.33,\;8.73)$ & 547\,882 & 2.08 \\
\midrule
Nugget $\gamma(0)$    & $\approx 21{,}283$ & 0.08 \\
Sill                  & 263\,003 & 1.000 \\
\textbf{Range} (80\,\% sill) & $h \approx 2.50$ & — \\
\bottomrule
\end{tabular}
\end{table}

The nugget is approximately 21\,000 (8\,\% of the sill), compared with
near-zero in Dataset~4.  This small but non-negligible nugget suggests
a modest level of measurement noise or micro-scale variability, implying
that RMSE cannot approach zero even with arbitrarily dense sampling.
The variogram range ($h \approx 2.50$) is similar to Dataset~4's 2.37.

\subsection{$k$-NN Oracle}

\begin{table}[h]
\centering
\caption{$k$-NN Oracle LOO-CV RMSE on the 40-sample train+val pool —
  Dataset~5 (original mg\,g$^{-1}$ space).}
\label{tab:ds5_knn_oracle}
\begin{tabular}{rr}
\toprule
$k$ & LOO RMSE (mg\,g$^{-1}$) \\
\midrule
1  & 328 \\
2  & 339 \\
3  & 310 \\
5  & 338 \\
7  & 380 \\
10 & 407 \\
\bottomrule
\end{tabular}
\end{table}

The best oracle ($k=3$) achieves RMSE $\approx 310$\,mg\,g$^{-1}$.
The GB model's 10-trial mean (193) is \emph{better} than the $k$-NN
oracle, which is expected: parametric models can extrapolate beyond the
nearest-neighbour boundary, while $k$-NN is purely interpolative.  The
LOO-CV RMSE of the GB model (248.8) is also better than the oracle,
confirming that the parametric model is extracting structural information
(not just local averaging).

\subsection{Nyquist Requirements}

Using $d_\text{eff} = 3.07$, $\bar{f} = 102.7$, and $L = 8.73$:

\begin{table}[h]
\centering
\caption{Nyquist sampling requirement — Dataset~5.}
\label{tab:ds5_nyquist_req}
\begin{tabular}{rrr}
\toprule
Target RMSE (mg\,g$^{-1}$) & Required spacing $\delta$ & $n_\text{required}$ \\
\midrule
300 & 5.84 &   3 \\
200 & 3.89 &  12 \\
150 & 2.92 &  29 \\
100 & 1.95 & 100 \\
 50 & 0.97 & 839 \\
\bottomrule
\end{tabular}
\end{table}

The Nyquist estimate for RMSE~$\leq 100$ requires $\sim$100 uniformly
distributed samples, consistent with the power-law projection of
$\approx 222$ training samples (the excess accounts for non-uniformity
of the current distribution).

\subsection{Data Augmentation Simulation}

\begin{table}[h]
\centering
\caption{$k$-NN ($k=3$) LOO RMSE on original 48 samples after adding
  interpolated synthetic training points — Dataset~5.}
\label{tab:ds5_augmentation}
\begin{tabular}{rr}
\toprule
$n_\text{extra}$ & LOO RMSE (mg\,g$^{-1}$) \\
\midrule
  0 & 288 \\
 10 & 288 \\
 20 & 284 \\
 30 & 247 \\
 50 & 172 \\
100 & 109 \\
\bottomrule
\end{tabular}
\end{table}

Adding 100 interpolated points reduces the $k$-NN oracle RMSE from
288 to 109\,mg\,g$^{-1}$ ($-$62\,\%), a steeper improvement than
Dataset~4's $-$22\,\% with the same augmentation count.

% -----------------------------------------------------------------------
\section{Dataset 5 — Recommendations}
\label{sec:ds5_recommendations}
% -----------------------------------------------------------------------

\subsection{Per-Decile Sampling Deficit}

\begin{table}[h]
\centering
\caption{Intra-decile density analysis — Dataset~5.
  Deficit estimated via
  $n_\text{extra} = n_\text{current}\,((\delta/\delta_\text{med})^{d_\text{eff}}-1)$,
  where $\delta_\text{med} = 0.882$.}
\label{tab:ds5_decile_deficit}
\small
\begin{tabular}{lrrr}
\toprule
Target decile (mg\,g$^{-1}$) & $n$ & Spacing & Deficit \\
\midrule
$[\phantom{00}13,\;\phantom{0}42)$ & 5 & 1.67 & $+30$ \\
$[\phantom{00}42,\;\phantom{0}90)$ & 5 & 1.41 & $+16$ \\
$[\phantom{00}90,\;194)$            & 5 & 1.12 & $+5$  \\
$[194,\;223)$                       & 4 & 3.31 & $+227$ \\
$[223,\;277)$                       & 5 & 1.54 & $+22$ \\
$[277,\;303)$                       & 5 & 1.71 & $+33$ \\
$[303,\;398)$                       & 4 & 1.44 & $+14$ \\
$[398,\;595)$                       & 5 & 0.99 & $+2$  \\
$[595,\;936)$                       & 5 & 1.71 & $+32$ \\
$[936,\;2251]$                      & 5 & 1.85 & $+43$ \\
\midrule
\textbf{Total deficit}              & & & $+424$ \\
\bottomrule
\end{tabular}
\end{table}

The most under-sampled decile ($[194, 223)$, deficit $+227$) is a narrow
range with only 4 samples spread over 3.31 spacing units — the highest
ratio in either dataset.  This appears to be a transition region between
low-performance and mid-performance adsorbents.

\subsection{Revised Data Collection Strategy}

\begin{table}[h]
\centering
\caption{Recommended data collection strategy — Dataset~5.}
\label{tab:ds5_recommendations}
\small
\begin{tabular}{p{4.5cm}rp{7cm}}
\toprule
Region & $n_\text{min}$ & Rationale \\
\midrule
SSA $\in [600,1400]$, dosage $\in [0.3,1.5]$
(transition zone) &
20 &
Bridges the diagonal clusters; fills the high-gradient
boundary. \\
\addlinespace
Target $y \in [194,223]$ and $[303,398]$
(thin deciles) &
15 &
These 4-sample deciles have the worst spacing and drive
interpolation error. \\
\addlinespace
Target $y \in [595,936]$ and $[936,2251]$
(high-performance tail) &
15 &
Total deficit = +75; high-value region most relevant to
applications. \\
\addlinespace
Neighbours of sample~43 and sample~5
(isolated extremes) &
5 &
Reduce NN-void contribution to extrapolation error. \\
\midrule
\textbf{Total minimum} & \textbf{55} & \\
\bottomrule
\end{tabular}
\end{table}

\subsection{Phased Roadmap}

\begin{table}[h]
\centering
\caption{Phased data collection roadmap — Dataset~5.}
\label{tab:ds5_roadmap}
\begin{tabular}{rlrr}
\toprule
Phase & Action & Additional $n$ & Projected RMSE \\
\midrule
Current & Baseline dataset & 48 total & $193 \pm 73$ \\
Phase 1 & +20 transition zone, +10 thin deciles & $+30$ & $\approx 150$--170 \\
Phase 2 & +15 high-tail, +5 isolated neighbours  & $+20$ & $\approx 120$--140 \\
Phase 3 & Stratified augmentation                & $+30$ & $\approx 90$--110 \\
\midrule
\textbf{Phase 3 total} & & \textbf{$\approx 128$} & \textbf{$\approx 100$} \\
\bottomrule
\end{tabular}
\end{table}

% -----------------------------------------------------------------------
\section{Dataset 5 — Summary and Comparison with Dataset 4}
\label{sec:ds5_conclusions}
% -----------------------------------------------------------------------

\begin{table}[h]
\centering
\caption{Head-to-head comparison of Dataset~4 and Dataset~5.}
\label{tab:ds5_comparison}
\small
\begin{tabular}{lrr}
\toprule
Metric & Dataset~4 & Dataset~5 \\
\midrule
$N$ (samples)                        & 56    & 48 \\
$d$ (features)                       & 7     & 6 \\
$d_\text{eff}$ (intrinsic dim.)      & 2.9   & 3.07 \\
Target mean $\pm$ std                & $479 \pm 542$ & $436 \pm 518$ \\
Target skewness                      & 2.13  & 2.27 \\
\midrule
10-trial mean RMSE                   & 279.9 & 193.0 \\
10-trial RMSE std                    & 111.9 & 72.8 \\
LOO-CV RMSE (original space)         & 345.9 & 248.8 \\
LOO-CV $R^2$                         & 0.571 & 0.794 \\
$k$-NN oracle RMSE                   & 328   & 310 \\
\midrule
Power-law exponent $b$               & $-0.166$ & $-0.465$ \\
Variogram range (80\,\% sill)        & 2.37  & 2.50 \\
Variogram nugget                     & $\approx 0$ & $\approx 21{,}000$ \\
$\varepsilon$ for full connectivity  & 7.30  & 4.00 \\
Max NN distance                      & 7.24  & 3.42 \\
\midrule
SSA importance                       & 0.203 & 0.384 \\
Feature 1 Spearman $\rho$ (SSA)      & 0.766 & 0.820 \\
\bottomrule
\end{tabular}
\end{table}

\paragraph{Key findings.}
\begin{enumerate}[leftmargin=*]
  \item \textbf{Dataset~5 yields markedly better prediction performance}
        ($-$31\,\% RMSE, $+$39\,\% $R^2$ on LOO-CV) despite having
        8 fewer samples and 1 fewer feature.  This suggests that the
        ``Equilibrium time'' variable in Dataset~4 introduced noise
        or redundancy rather than useful signal.

  \item \textbf{The learning curve is no longer saturated.}  The steeper
        power-law exponent ($-0.47$ vs.\ $-0.17$) means additional data
        will produce substantial RMSE reductions.  Roughly 50 training
        samples are needed for RMSE~$\leq 200$, compared with
        $\approx 1{,}124$ in Dataset~4.

  \item \textbf{Feature-space topology is more compact.}  Full connectivity
        at $\varepsilon = 4.0$ (vs.\ 7.3) and a maximum NN distance of
        3.42 (vs.\ 7.24) indicate fewer extreme voids, enabling more
        reliable interpolation.

  \item \textbf{SSA is the single most important feature} (importance 0.38,
        Spearman $\rho = 0.820$), followed by Dosage ratio (0.30).
        Together they account for 68\,\% of model importance.

  \item \textbf{A non-zero variogram nugget ($\approx 21{,}000$)} appeared
        in Dataset~5, suggesting modest measurement noise that places
        a fundamental lower bound on achievable RMSE of approximately
        $\sqrt{21{,}000} \approx 146$\,mg\,g$^{-1}$ — or lower if the
        nugget is partly attributable to short-range spatial variability.

  \item \textbf{Recommended minimum additional data: 55 targeted experiments}
        (Phase~1+2), with a goal of reaching RMSE~$\approx 100$\,mg\,g$^{-1}$
        at 128 total samples (Phase~3).
\end{enumerate}

% -----------------------------------------------------------------------
% Model configuration for Dataset 5
% -----------------------------------------------------------------------
\section*{Appendix D: Dataset~5 Statistical Tests Summary}
\label{app:ds5_stats}
\addcontentsline{toc}{section}{Appendix D: Dataset~5 Statistical Tests Summary}

\begin{table}[h]
\centering
\caption{Key statistical results — Dataset~5.}
\label{tab:ds5_stats_summary}
\begin{tabular}{lll}
\toprule
Metric & Value & Interpretation \\
\midrule
Test R$^2$ mean $\pm$ std & $0.638 \pm 0.535$ & Very high variance \\
Test RMSE mean $\pm$ std & $193 \pm 73$ & CV = 0.38 \\
LOO-CV RMSE (original space) & 248.8 & Near-unbiased, single estimate \\
LOO-CV R$^2$ (original space) & 0.794 & Substantially better than 10-trial mean \\
95\,\% Bootstrap CI (RMSE) & (149, 235) & $\pm$22\,\% of mean \\
Test set Jaccard similarity & $0.115 \pm 0.091$ & Near-independent folds \\
Pearson(RMSE, test mean) & $r = 0.375$ & Not significant ($p=0.285$) \\
Pearson(R$^2$, test max) & $r = 0.622$ & Borderline ($p=0.055$) \\
Power-law exponent & $-0.465$ & Near root-$n$ learning \\
Power-law fit $r$ & $-0.957$ & Good fit \\
\midrule
Spearman(SSA, target) & $\rho = 0.820$ & Strongest feature \\
Spearman(Dosage, target) & $\rho = 0.571$ & Second strongest \\
Spearman(PV, target) & $\rho = 0.651$ & Third strongest \\
\bottomrule
\end{tabular}
\end{table}

The model configuration is identical to Dataset~4 (Appendix~\ref{app:model}):
GradientBoostingRegressor with the same hyperparameters.

% -----------------------------------------------------------------------
\section*{Appendix C: Topological \& Nyquist Summary Statistics}
\label{app:topo}
\addcontentsline{toc}{section}{Appendix C: Topological \& Nyquist Summary Statistics}
% -----------------------------------------------------------------------

\begin{table}[h]
\centering
\caption{Key topology and Nyquist analysis results.}
\label{tab:topo_summary}
\begin{tabular}{lll}
\toprule
Quantity & Value & Interpretation \\
\midrule
Ambient dimensionality $d$ & 7 & Measured features \\
Intrinsic dimensionality $d_\text{eff}$ & 2.9 & 2NN estimate \\
Correlation dimension & 2.18 & Power-law fit on pairwise dist \\
PCA 90\,\% variance & 5 PCs & Remaining 2 PCs carry $<$10\,\% \\
\midrule
Median NN distance & 0.806 & Normalised feature space \\
Maximum NN distance & 7.24 & Sample~30 is an extremal outlier \\
$\varepsilon$ for full connectivity & 7.30 & vs.\ median 0.81; factor 9 \\
Isolated samples ($>2\times$median) & 7 & 12.5\,\% of dataset \\
\midrule
Variogram range (80\,\% sill) & 2.37 & Correlation length of target \\
Variogram nugget & $\approx 0$ & No irreducible noise \\
Median spatial frequency $\bar{f}$ & 139 mg\,g$^{-1}$/unit & \\
Max spatial frequency $f_\text{max}$ & 1\,708 mg\,g$^{-1}$/unit & \\
\midrule
$k$-NN oracle RMSE (best $k$) & 328 mg\,g$^{-1}$ & Interpolation ceiling (1-NN LOO) \\
Best GB model RMSE (10-trial) & 280 mg\,g$^{-1}$ & 85\,\% of oracle \\
Best GB model RMSE (LOO-CV) & 346 mg\,g$^{-1}$ & Near-oracle, same units \\
Augmented oracle (+100 pts) & 285 mg\,g$^{-1}$ & Validates Nyquist prediction \\
\midrule
Nyquist $n$ for RMSE 100 & $\sim$63 uniform & Idealised estimate \\
Practical $n$ for RMSE 100 & $\sim$130 targeted & Accounting for non-uniformity \\
Total sample deficit (decile) & 636 & For uniform median spacing \\
\bottomrule
\end{tabular}
\end{table}

% -----------------------------------------------------------------------
\bibliographystyle{plainnat}
\begin{thebibliography}{99}

\bibitem[Arlot \& Celisse(2010)]{Arlot2010}
Arlot, S. \& Celisse, A. (2010).
\newblock A survey of cross-validation procedures for model selection.
\newblock \emph{Statistics Surveys}, 4, 40--79.
\newblock \doi{10.1214/09-SS054}

\bibitem[Balestriero et al.(2021)]{Balestriero2021}
Balestriero, R., Pesenti, J., \& LeCun, Y. (2021).
\newblock Learning in high dimension always amounts to extrapolation.
\newblock \emph{arXiv preprint arXiv:2110.09485}.

\bibitem[Bellman(1957)]{Bellman1957}
Bellman, R. (1957).
\newblock \emph{Dynamic Programming}.
\newblock Princeton University Press, Princeton, NJ.

\bibitem[Carlsson(2009)]{Carlsson2009}
Carlsson, G. (2009).
\newblock Topology and data.
\newblock \emph{Bulletin of the American Mathematical Society}, 46(2), 255--308.
\newblock \doi{10.1090/S0273-0979-09-01249-X}

\bibitem[Cover \& Hart(1967)]{Cover1967}
Cover, T. M. \& Hart, P. E. (1967).
\newblock Nearest neighbor pattern classification.
\newblock \emph{IEEE Transactions on Information Theory}, 13(1), 21--27.
\newblock \doi{10.1109/TIT.1967.1053964}

\bibitem[Cressie(1993)]{Cressie1993}
Cressie, N. A. C. (1993).
\newblock \emph{Statistics for Spatial Data} (revised edition).
\newblock Wiley, New York.

\bibitem[Edelsbrunner \& Harer(2010)]{Edelsbrunner2010}
Edelsbrunner, H. \& Harer, J. L. (2010).
\newblock \emph{Computational Topology: An Introduction}.
\newblock American Mathematical Society, Providence, RI.

\bibitem[Facco et al.(2017)]{Facco2017}
Facco, E., d'Errico, M., Rodriguez, A., \& Laio, A. (2017).
\newblock Estimating the intrinsic dimension of datasets by a minimal
  neighborhood information.
\newblock \emph{Scientific Reports}, 7, 12140.
\newblock \doi{10.1038/s41598-017-11873-y}

\bibitem[Grassberger \& Procaccia(1983)]{Grassberger1983}
Grassberger, P. \& Procaccia, I. (1983).
\newblock Measuring the strangeness of strange attractors.
\newblock \emph{Physica D: Nonlinear Phenomena}, 9(1--2), 189--208.
\newblock \doi{10.1016/0167-2789(83)90298-1}

\bibitem[Hastie et al.(2009)]{Hastie2009}
Hastie, T., Tibshirani, R., \& Friedman, J. (2009).
\newblock \emph{The Elements of Statistical Learning: Data Mining, Inference,
  and Prediction} (2nd ed.).
\newblock Springer, New York.
\newblock \doi{10.1007/978-0-387-84858-7}

\bibitem[Jolliffe(2002)]{Jolliffe2002}
Jolliffe, I. T. (2002).
\newblock \emph{Principal Component Analysis} (2nd ed.).
\newblock Springer, New York.
\newblock \doi{10.1007/b98835}

\bibitem[Matheron(1963)]{Matheron1963}
Matheron, G. (1963).
\newblock Principles of geostatistics.
\newblock \emph{Economic Geology}, 58(8), 1246--1266.
\newblock \doi{10.2113/gsecongeo.58.8.1246}

\bibitem[McKay et al.(1979)]{McKay1979}
McKay, M. D., Beckman, R. J., \& Conover, W. J. (1979).
\newblock A comparison of three methods for selecting values of input variables
  in the analysis of output from a computer code.
\newblock \emph{Technometrics}, 21(2), 239--245.
\newblock \doi{10.1080/00401706.1979.10489755}

\bibitem[Nyquist(1928)]{Nyquist1928}
Nyquist, H. (1928).
\newblock Certain topics in telegraph transmission theory.
\newblock \emph{Transactions of the American Institute of Electrical
  Engineers}, 47(2), 617--644.
\newblock \doi{10.1109/T-AIEE.1928.5055024}

\bibitem[Pronzato \& M{\"u}ller(2012)]{Pronzato2012}
Pronzato, L. \& M{\"u}ller, W. G. (2012).
\newblock Design of computer experiments: space filling and beyond.
\newblock \emph{Statistics and Computing}, 22(3), 681--701.
\newblock \doi{10.1007/s11222-011-9242-3}

\bibitem[Settles(2012)]{Settles2012}
Settles, B. (2012).
\newblock \emph{Active Learning}.
\newblock Synthesis Lectures on Artificial Intelligence and Machine Learning.
\newblock Morgan \& Claypool, San Rafael, CA.
\newblock \doi{10.2200/S00429ED1V01Y201207AIM018}

\bibitem[Shannon(1949)]{Shannon1949}
Shannon, C. E. (1949).
\newblock Communication in the presence of noise.
\newblock \emph{Proceedings of the IRE}, 37(1), 10--21.
\newblock \doi{10.1109/JRPROC.1949.232969}

\bibitem[Stone(1977)]{Stone1977}
Stone, C. J. (1977).
\newblock Consistent nonparametric regression.
\newblock \emph{The Annals of Statistics}, 5(4), 595--620.
\newblock \doi{10.1214/aos/1176343886}

\bibitem[Webster \& Oliver(2007)]{Webster2007}
Webster, R. \& Oliver, M. A. (2007).
\newblock \emph{Geostatistics for Environmental Scientists} (2nd ed.).
\newblock Wiley, Chichester.
\newblock \doi{10.1002/9780470517277}

\bibitem[Zomorodian \& Carlsson(2005)]{Zomorodian2005}
Zomorodian, A. \& Carlsson, G. (2005).
\newblock Computing persistent homology.
\newblock \emph{Discrete \& Computational Geometry}, 33(2), 249--274.
\newblock \doi{10.1007/s00454-004-1146-y}

\end{thebibliography}

\end{document}

